% GENERATED FILE. Edit canonical lessons, then run npm run generate:books.
% canonical-source-hash: fnv1a64:15e978266f8938cf
% canonical-lessons: ZH-C34-jiao, ZH-W34-jiao, ZH-C34-nin, ZH-C34-xing, ZH-C34-renshi, ZH-C34-pengyou

\chapter{What You Are Called}
\label{ch:zh-what-you-are-called}
\hlchaptermodality{34}

\begin{chapteropening}
\textbf{By the end of this chapter:} \emph{I can ask somebody's name and family name, politely too.}

Complete the last lesson of chapter 34: I can ask somebody's name and family name, politely too.
\end{chapteropening}

\section[\texorpdfstring{\zh{叫} (jiào)}{叫 (jiào)}]{\zh{叫} (jiào) — to be called; to call}

\label{lesson:ZH-C34-jiao}



\begin{quote}
Before the new one: say the Mandarin for younger sister, then the Mandarin for child.
\end{quote}

\subsection*{You'll want to know: \zh{叫}}
\textbf{\zh{叫}} — \emph{jiào} — \textquotedblleft{}to be called; to call\textquotedblright{}.

\begin{quote}
\textbf{\zh{你叫什么名字}} \emph{nǐ jiào shénme míngzi} — What's your name?
\end{quote}

The new character, in pieces: \zh{叫} — \zh{口}, the mouth piece, small on the left, and on the right two strokes: one that flicks up, and a long vertical. The next lesson puts it in your hand.

\begin{cousinweb}
\textbf{\zh{叫}} (\emph{jiào}) is inherited from Middle Chinese \emph{kewH} (Guangyun, in Baxter's transcription), glossed 'to call out'. The sense 'to be named' is not given there; from calling a person to what a person is called is a likely step, but it is not settled.
\end{cousinweb}

\begin{grammarlens}[title={What you've built}]
One more word about people.
\end{grammarlens}

\subsection*{Your turn}
\begin{itemize}
\raggedright
  \item \emph{Say it:} \emph{jiào}
  \item \emph{Say it:} \emph{jiào}, once more
  \item \emph{Say it:} \emph{nǐ jiào shénme míngzi} — What's your name?
  \item \emph{Recall:} say the Mandarin for younger sister, then the Mandarin for child, then say \emph{jiào} again
\end{itemize}

\begin{tcolorbox}[breakable,colback=teal!4,colframe=teal!35!black,title={Before you move on}]
What does \zh{叫} mean? (\textquotedblleft{}To be called; to call\textquotedblright{}.) Say it once more.
\end{tcolorbox}
\section[\texorpdfstring{\zh{叫} (jiào)}{叫 (jiào)}]{\zh{叫} — a small mouth, then two strokes}

\label{lesson:ZH-W34-jiao}



\begin{quote}
The word first, then the character inside it.

\begin{itemize}
\raggedright
  \item \emph{Recall:} say \emph{jiào} — \textbf{R1}, one lesson back, before it decays
\end{itemize}

\begin{quote}
\zh{口} \zh{吗}
\end{quote}

\emph{Read it:} each one aloud, with nothing said first — \emph{kǒu}, \emph{ma}
\end{quote}

\subsection*{Script — a small mouth, then two strokes}
\hlblockfigure{\detokenize{figures/ZH-W34-jiao-filmstrip.pdf}}{How \zh{叫} is written, stroke by stroke}

\begin{quote}
\zh{叫}
\end{quote}

\textbf{Two pieces, left before right.}

\textbf{The left piece, \zh{口}, three strokes}, as in \zh{吗}, small and set a little high: the left side, the top that turns down the right, the bottom that closes it.

\textbf{The right piece, two strokes.} A vertical that flicks up to the right at the foot. Then a long vertical, the tallest stroke in the character.

3 + 2 = five.

\subsection*{Your turn}
\begin{itemize}
\raggedright
  \item \emph{Write it:} \zh{叫}, all five, with the model in front of you
  \item \emph{Cover:} write it again from memory, and count the pieces, not the strokes
\end{itemize}

\begin{tcolorbox}[breakable,colback=teal!4,colframe=teal!35!black,title={Before you move on}]
How many pieces in \zh{叫}? (\textbf{Two} — \zh{口}, then two strokes.) Which character you know has the same small \zh{口} on the left? (\textbf{\zh{吗}}.) Which stroke is last? (\textbf{The long vertical}.)

Source: \href{https://raw.githubusercontent.com/chanind/hanzi-writer-data/68d10a4b21150cae5e1ebbd223eed289cf32d90c/data/\%E5\%8F\%AB.json}{Hanzi Writer Data \zh{叫}, pinned stroke paths and medians}.
\end{tcolorbox}
\section[nín]{nín — you (polite)}

\label{lesson:ZH-C34-nin}



\begin{quote}
Before the new one: say the Mandarin for child, then the Mandarin for to be called.
\end{quote}

\subsection*{You'll want to know: nín}
\emph{nín} — \textquotedblleft{}you (polite)\textquotedblright{}.

\begin{quote}
\emph{nín jiào shénme míngzi} — May I ask your name?
\end{quote}

Its characters come later; for now it is a word for the ear.

\begin{cousinweb}
\emph{nín} is the polite form of \zh{你}, for older people and strangers. No dictionary we use records an older form, so its history is not settled.
\end{cousinweb}

\begin{grammarlens}[title={What you've built}]
One more word about people.
\end{grammarlens}

\subsection*{Your turn}
\begin{itemize}
\raggedright
  \item \emph{Say it:} \emph{nín}
  \item \emph{Say it:} \emph{nín}, once more
  \item \emph{Say it:} \emph{nín jiào shénme míngzi} — May I ask your name?
  \item \emph{Recall:} say the Mandarin for child, then the Mandarin for to be called, then say \emph{nín} again
\end{itemize}

\begin{tcolorbox}[breakable,colback=teal!4,colframe=teal!35!black,title={Before you move on}]
What does nín mean? (\textquotedblleft{}You (polite)\textquotedblright{}.) Say it once more.
\end{tcolorbox}
\section[xìng]{xìng — to have the family name; family name}

\label{lesson:ZH-C34-xing}



\begin{quote}
Before the new one: say the Mandarin for to be called, then the Mandarin for you (polite).
\end{quote}

\subsection*{You'll want to know: xìng}
\emph{xìng} — \textquotedblleft{}to have the family name; family name\textquotedblright{}.

\begin{quote}
\emph{nǐ xìng shénme} — What's your family name?
\end{quote}

Its characters come later; for now it is a word for the ear.

\begin{cousinweb}
\emph{xìng} is inherited from Middle Chinese \emph{sjengH} (Guangyun, in Baxter's transcription), glossed 'clan name'.
\end{cousinweb}

\begin{grammarlens}[title={What you've built}]
One more word about people.
\end{grammarlens}

\subsection*{Your turn}
\begin{itemize}
\raggedright
  \item \emph{Say it:} \emph{xìng}
  \item \emph{Say it:} \emph{xìng}, once more
  \item \emph{Say it:} \emph{nǐ xìng shénme} — What's your family name?
  \item \emph{Recall:} say the Mandarin for to be called, then the Mandarin for you (polite), then say \emph{xìng} again
\end{itemize}

\begin{tcolorbox}[breakable,colback=teal!4,colframe=teal!35!black,title={Before you move on}]
What does xìng mean? (\textquotedblleft{}To have the family name; family name\textquotedblright{}.) Say it once more.
\end{tcolorbox}
\section[rènshi]{rènshi — to know someone; to recognise}

\label{lesson:ZH-C34-renshi}



\begin{quote}
Before the new one: say the Mandarin for you (polite), then the Mandarin for to have the family name.
\end{quote}

\subsection*{You'll want to know: rènshi}
\emph{rènshi} — \textquotedblleft{}to know someone; to recognise\textquotedblright{}.

\begin{quote}
\emph{wǒ rènshi tā} — I know him.
\end{quote}

Its characters come later; for now it is a word for the ear.

\begin{cousinweb}
\emph{rènshi} is built from two old words for knowing, both inherited from Middle Chinese: \emph{nyinH}, which the Guangyun glosses with the second, and \emph{syik}, 'to know' (Guangyun, in Baxter's transcription).
\end{cousinweb}

\begin{grammarlens}[title={What you've built}]
One more word about people.
\end{grammarlens}

\subsection*{Your turn}
\begin{itemize}
\raggedright
  \item \emph{Say it:} \emph{rènshi}
  \item \emph{Say it:} \emph{rènshi}, once more
  \item \emph{Say it:} \emph{wǒ rènshi tā} — I know him.
  \item \emph{Recall:} say the Mandarin for you (polite), then the Mandarin for to have the family name, then say \emph{rènshi} again
\end{itemize}

\begin{tcolorbox}[breakable,colback=teal!4,colframe=teal!35!black,title={Before you move on}]
What does rènshi mean? (\textquotedblleft{}To know someone; to recognise\textquotedblright{}.) Say it once more.
\end{tcolorbox}
\section[péngyou]{péngyou — friend}

\label{lesson:ZH-C34-pengyou}



\begin{quote}
Before the new one: say the Mandarin for to have the family name, then the Mandarin for to know someone.
\end{quote}

\subsection*{You'll want to know: péngyou}
\emph{péngyou} — \textquotedblleft{}friend\textquotedblright{}.

\begin{quote}
\emph{tā shì wǒ de péngyou} — He's a friend of mine.
\end{quote}

Its characters come later; for now it is a word for the ear.

\begin{cousinweb}
\emph{péngyou} is built from two words for companions, both inherited from Middle Chinese: \emph{bong} 'a band of companions' and \emph{hjuwX} 'friend' (Guangyun, in Baxter's transcription), which the dictionary explains as people of one mind.
\end{cousinweb}

\begin{grammarlens}[title={What you've built}]
One more word about people.
\end{grammarlens}

\subsection*{Your turn}
\begin{itemize}
\raggedright
  \item \emph{Say it:} \emph{péngyou}
  \item \emph{Say it:} \emph{péngyou}, once more
  \item \emph{Say it:} \emph{tā shì wǒ de péngyou} — He's a friend of mine.
  \item \emph{Recall:} say the Mandarin for to have the family name, then the Mandarin for to know someone, then say \emph{péngyou} again
\end{itemize}

\begin{tcolorbox}[breakable,colback=teal!4,colframe=teal!35!black,title={Before you move on}]
What does péngyou mean? (\textquotedblleft{}Friend\textquotedblright{}.) Say it once more.
\end{tcolorbox}
